% The modelled COMEX-London gold spread, as a system of equations.
%
% Typeset version of section 9 of claude/spread-explainer/SPREAD_EXPLAINED.html
% and section 4 of spread_model_note.pdf, where the same system is set in
% monospace. Equation numbers match those documents.
%
% Build (from the repo root), keeping aux files out of the working tree:
%   pdflatex -interaction=nonstopmode -output-directory=<build> \
%     claude/spread-explainer/spread_system.tex          % run it twice
%
% Use pdflatex directly rather than latexmk: MiKTeX 23.5 on this machine
% prints "You are running MiKTeX on an unsupported version of Windows"
% (Windows 11 build 26200) and exits non-zero even when the compile
% succeeded, which latexmk treats as fatal and refuses to work through.
%
% Package set is deliberately the same one the repo's other .tex already
% compiles with, so no package has to be fetched at build time.

\documentclass[11pt,a4paper]{article}

\usepackage[margin=1in]{geometry}
\usepackage{amsmath}
\usepackage{amssymb}
\usepackage{url}
\usepackage{booktabs}
\usepackage{enumitem}
\usepackage{parskip}
\usepackage[hidelinks]{hyperref}

% The objects of the model, named once so the notation cannot drift.
\newcommand{\sN}{s^{N}}                        % New York storage cost
\newcommand{\lN}{\ell^{N}_{t}}                 % New York lease rate
\newcommand{\lL}{\ell^{L}_{t}}                 % London lease rate
\newcommand{\kLN}{\kappa^{L\to N}_{t}}         % cost of moving London -> New York
\newcommand{\kNL}{\kappa^{N\to L}_{t}}         % and back
\newcommand{\Fhat}{\widehat{F}}
\newcommand{\disloc}{\operatorname{disloc}}
\newcommand{\Var}{\operatorname{Var}}

\title{\textbf{The Modelled Spread}\\[0.4em]
\large A System of Equations for the COMEX--London Gold Basis}
\author{Methods note, GOLD project}
\date{16 September 2026}

\begin{document}
\maketitle

\noindent
The spread between the COMEX gold futures price and the London benchmark is
usually reported as a single annualised number. That number mixes three
economically distinct objects: a premium for metal being in New York rather
than London, a premium for the risk that a tariff lands before delivery, and
the cost of carrying metal through time. This note writes the system down so
the three can be held apart, states which combinations the data identify, and
records what the estimated versions of them look like.

Equation numbers match Section~9 of \path{SPREAD_EXPLAINED.html} and
Section~4 of \path{spread_model_note.pdf}. All rates are annualised
decimals, $\tau$ is in calendar days to first notice, and prices are US
dollars per troy ounce.

\section*{Notation}

\begin{center}
\begin{tabular}{@{}lll@{}}
\toprule
\textbf{Symbol} & \textbf{Meaning} & \textbf{Status} \\
\midrule
$S_t$ & London price, unallocated, 15:00 London & observed \\
$F_t(\tau)$ & COMEX price for delivery in $\tau$ days & observed, 13:30 New York \\
$V_t(\tau)$ & value today of metal deliverable in New York at $\tau$ & latent \\
$p_t$ & location premium, fraction of spot (place, form, trust) & latent; estimated \\
$r_t$ & dollar interest rate (SOFR; fed funds before Apr 2018) & observed \\
$\sN$ & New York storage cost, per year & assumed; see note below \\
$\ell^{N}_{t},\ \ell^{L}_{t}$ & lease rate on metal in New York, in London & latent \\
$h_t,\ \theta_t$ & hazard rate of a tariff; tariff rate if it lands & latent \\
$\Lambda_t(\tau)$ & probability a delivery at $\tau$ is tariffed & latent \\
$\kappa_t$ & all-in cost of moving an ounce between locations & latent; estimated \\
$\varepsilon_t$ & gold's return from 15:00 London to the settlement & latent; measurable \\
$Q_t$ & metal shipped, tonnes a month & observed (Swiss customs) \\
$\bar{Q}$ & refining and secure-freight capacity, tonnes a month & latent \\
\bottomrule
\end{tabular}
\end{center}

\section*{The system}

\noindent\textbf{Pricing through time, within New York.}
\begin{align}
F_t(\tau) &= V_t(\tau)\cdot\exp\!\left[\left(r_t+\sN-\lN\right)\frac{\tau}{365}\right]
\end{align}

\noindent\textbf{Pricing across locations, and the policy risk inside it.}
\begin{align}
V_t(\tau) &= S_t\left[\,1+p_t+\Lambda_t(\tau)\,\theta_t\,\right] \\[2pt]
\Lambda_t(\tau) &= 1-\exp\!\left(-h_t\frac{\tau}{365}\right)
   \;\approx\; h_t\frac{\tau}{365}-\frac{1}{2}\left(h_t\frac{\tau}{365}\right)^{\!2}
\end{align}

\noindent\textbf{What arbitrage enforces: a band, not a point.}
\begin{align}
-\kNL \;&\le\; p_t S_t \;\le\; \kLN
   \qquad\text{unless capacity binds} \\[2pt]
\kLN &= \underbrace{c_{\text{freight}}+c_{\text{ins}}+c_{\text{recast}}}_{\text{physical, per ounce}}
   \;+\; \underbrace{\left(r_t+\lL\right)\frac{d_{\text{transit}}}{365}\,S_t}_{\text{financing in transit}}
\end{align}

\noindent\textbf{Measurement: the two legs are not observed at the same instant.}
\begin{align}
\Fhat_t(\tau) &= F_t(\tau)\cdot e^{\varepsilon_t},
   \qquad \mathbb{E}[\varepsilon_t]=\mu,
   \qquad \Var[\varepsilon_t]=\sigma^{2}
\end{align}
The London benchmark is struck at 15:00 London and the COMEX settlement about
three and a half hours later, so $\varepsilon_t$ is gold's return over that
window. Every contract settles at the same instant, so $\varepsilon_t$ is
\emph{common} to all of them on date~$t$: it shifts the whole curve without
tilting it.

\noindent\textbf{Flow, with a capacity ceiling.}
\begin{align}
Q_t &= \min\!\left\{\bar{Q},\;\;\beta\max\!\left(0,\;p_t S_t-\kappa_t\right)\right\}+u_t
\end{align}

\section*{What is estimated}

Taking logs of (1) and (2), substituting (3) and (6), and collecting powers of
$\tau$ gives a curve that is quadratic in the horizon. This is the regression
run each day across that date's contracts, weighted by the square root of open
interest:
\begin{align}
\ln \Fhat_t(\tau_i) &= \alpha_t+\beta_t\tau_i+\gamma_t\tau_i^{2}+\eta_{it}
\end{align}
The three fitted coefficients map onto the objects above:
\begin{align}
\hat{p}_t \;\equiv\; \alpha_t-\ln S_t &= p_t+\varepsilon_t \\[2pt]
\hat{\omega}_t \;\equiv\; 365\beta_t-r_t &= \sN-\lN+h_t\theta_t \\[2pt]
\gamma_t &= -\,\frac{h_t^{2}\theta_t}{2\cdot 365^{2}} \;\le\; 0
\end{align}

\section*{The modelled spread}

\begin{align}
b_t(\tau)\;\equiv\;F_t(\tau)-S_t
  \;&=\; S_t\left[\;
  \underbrace{p_t}_{\substack{\text{place, form,}\\ \text{trust}}}
  \;+\;\underbrace{\Lambda_t(\tau)\,\theta_t}_{\text{policy risk}}
  \;+\;\underbrace{\left(r_t+\sN-\lN\right)\dfrac{\tau}{365}}_{\text{carry}}
  \;\right] \\[6pt]
\disloc_t(T) \;&\approx\; p_t\,\frac{365}{T}\;+\;h_t\theta_t\;+\;\left(\sN-s\right)\;-\;\lN
\end{align}

Equation (12) is the spread itself. Equation (13) is what the project's
headline number does to it: reading the fitted curve at a fixed horizon $T$
and subtracting an assumed carry rate $s$ divides the location premium by the
horizon, leaves the policy and lease terms alone, and adds whatever error sits
in the storage assumption. The multiplier $365/T$ is $12.2$ at thirty days,
$4.06$ at ninety and $2.03$ at one hundred and eighty, so the horizon silently
sets how much of the reported number is location and how much is carry. That
is the case for reporting $\hat{p}_t$ and $\hat{\omega}_t$ separately rather
than any single dislocation.

\section*{What the system implies, and what the data say}

Estimated over 2{,}857 daily fits, 2015--2026 (\texttt{spread\_model\_checks.py}):

\begin{itemize}[leftmargin=1.2em,itemsep=4pt]
\item \textbf{The timing error lands in the level, not the slope.} Because
$\varepsilon_t$ is common across contracts it enters (9) and not (10).
Regressing $\hat{p}_t$ on the day's London-to-London return gives
$+0.193$ (s.e.\ $0.015$), against $3.5/24=0.146$ predicted by a random walk
over the window; the same regression for $\hat{\omega}_t$ gives $-0.012$
(s.e.\ $0.011$), indistinguishable from zero. Any result resting on the slope
needs no timing correction.

\item \textbf{The band in (4) is where the transfer cost is.} In the calm
years 2015--19 the fitted level averages $-\$0.55$ an ounce, so arbitrage does
hold $p_t S_t$ near zero. Trade estimates of the all-in cost of moving London
metal into COMEX form are roughly \$1.5--2.5 an ounce, and the flow threshold
estimated from (7) is \$3.77 with a 90\% interval of \$0.83--10.25. Three
independent readings of $\kappa_t$, mutually consistent.

\item \textbf{(10) does not identify the lease rate on its own.} The calm-period
wedge is $\hat{\omega}=+0.58$ points a year. With $\ell^{N}\approx 0$ in calm
years that implies storage of about $0.58\%$, not the $0.25\%$ the code
assumes; if storage really is $0.25\%$, the implied lease rate is $-0.33\%$,
which nobody quotes. One of the two is mis-set, and because $s$ enters (13)
one-for-one, so is every reported dislocation level. Levels should be read in
differences until this is pinned down.

\item \textbf{Curvature is the lever on policy risk.} $\gamma_t$ is negative on
average, with mean over standard error of $-10.7$ in calm years and $-22.0$ in
the tariff window --- a roughly sevenfold deepening. Consistent with the
hazard term in (3), but present in calm years too, so the dated resolutions
(the April 2025 exemption, the July 2025 customs ruling) are still needed to
attribute it: at resolution $h_t\to 0$, so the slope and curvature should step
while the level does not.

\item \textbf{The flow equation is not a response function.} Shipments respond
to the premium and also close it, so estimating (7) on monthly data recovers
an equilibrium relation rather than the supply of arbitrage. The coefficient
is descriptive; identification needs the timing within the month, or an event
study around the dated announcements.
\end{itemize}

\vspace{0.5em}
\noindent\rule{\linewidth}{0.4pt}\\[0.3em]
{\footnotesize Sources: \path{claude/spread-explainer/spread_model_checks.py}
and \path{worked_examples.py}, both reading \path{data/processed/}; the
kink threshold is from \path{claude/mechanism-figures/estimate_kink.py}.
Where this note disagrees with earlier project text, it says so.}

\end{document}
